\documentclass[10pt,a4paper]{amsart}
\usepackage[margin=19mm]{geometry}
\usepackage{amsmath,amssymb,mathtools}
\usepackage[scr=rsfs,cal=euler]{mathalfa}
\usepackage{xfrac}
\usepackage[unicode,hidelinks,bookmarks=false]{hyperref}
\newcommand{\ie}{i.e.\ }
\newcommand{\cf}{cf.\ }
\newcommand{\resp}{respectively\ }
\newcommand{\Subsection}{\subsection}
\providecommand{\nobreakdash}{\nobreak\hbox{-}}
\setlength{\emergencystretch}{2em}
\multlinegap0pt
\usepackage{amssymb}
\newtheorem{theorem}{Theorem}
\title{A Simple Trigonometric Classification of Quartic Roots}
\author{Sawon Pratiher}
\date{Source: arXiv:2603.28814v1 (CC BY 4.0)}
\begin{document}
\maketitle

\noindent\emph{Source and licence.} A Simple Trigonometric Classification of Quartic Roots. Version arXiv:2603.28814v1 (CC BY 4.0), distributed by arXiv under the Creative Commons Attribution 4.0 International licence, http://creativecommons.org/licenses/by/4.0/, as recorded in the arXiv licence metadata for this exact version and shown on the abstract page https://arxiv.org/abs/2603.28814v1. Full licence notice: \texttt{licenses/arxiv-2603.28814-CC-BY-4.0.txt}.\par\smallskip

\noindent\textit{Editorial scope.} Counted unit: the article's theorem thm:four\_complex (source0.tex lines 165--170). The article's complete original proof of it is delivered with it (source0.tex lines 172--178, one \texttt{proof} environment of 7 lines). No mathematics is written, repaired or re-derived: the statement and the proof are the article's own lines, in the article's own notation, and no retained line is substituted or re-flowed.  Retained uncounted as the source-local context the proof needs: lines 84--86.  Nothing was omitted from any retained span, so the package carries no omission record.  No retained line contains a citation, so the wrapper carries no bibliography and no bibliography entry.  No retained line contains a figure environment, an image-inclusion command or drawing code, so no image is delivered and the package declares no asset dependency.  Editorial formatting only, in addition: an AMS \texttt{amsart} preamble that restates the article's own declarations and macros used by the retained lines, namely the environments \texttt{theorem} on separate counters, together with the standard packages those lines need.  The retained environments print in this document as Theorem 1 in order of appearance, whereas the article prints the same items with the numbering of its own full document.  The complete native source of the article as distributed in the arXiv e-print for this exact version is bundled unchanged as \texttt{sources/197469/source0.tex}.
\begin{equation}\label{eq:boundary}
f(0) = \frac{8\,P(u)}{u^4}, \qquad f(\pi) = \frac{8\,P(-u)}{u^4}.
\end{equation}

\begin{theorem}[All four roots complex]\label{thm:four_complex}
Suppose $m < 0$. All four roots of $P(t) = 0$ are complex if and only if $f(\theta) > 0$ for all $\theta \in [0,\pi]$ (equivalently, $N_{\mathrm{int}} = 0$ and $N_{\mathrm{ext}} = 0$). A sufficient condition is
\begin{equation}\label{eq:four_complex_suff}
b > |a| + 1.
\end{equation}
\end{theorem}

\begin{proof}
If $f > 0$ on $[0,\pi]$, then $N_{\mathrm{int}} = 0$. The boundary values $f(0)$ and $f(\pi)$ are both positive, so $P(u) > 0$ and $P(-u) > 0$ by (\ref{eq:boundary}), giving $N_{\mathrm{ext}} = 0$. Thus $N_{\mathrm{real}} = 0$.

For the sufficient condition: for all $\theta$, $|a\cos\theta + \cos 4\theta| \leq |a| + 1$, so $f(\theta) \geq b - (|a|+1) > 0$ when $b > |a| + 1$.

Conversely, if $f < 0$ on $[0,\pi]$, then $N_{\mathrm{int}} = 0$ but $f(0) < 0$ and $f(\pi) < 0$, giving $N_{\mathrm{ext}} = 2$, so there are two (not zero) real roots. Thus the only route to zero real roots is $f > 0$.
\end{proof}

\end{document}
